% GENERATED FILE. Edit canonical lessons, then run npm run generate:books.
% canonical-source-hash: fnv1a64:5980f467d59a94d8
% canonical-lessons: RU-C293-sud, RU-C293-advokat, RU-C293-prigovor, RU-C293-tyurma, RU-C293-obvinyat

\chapter{In court}
\label{ch:ru-in-court}
\hlchaptermodality{293}

\begin{chapteropening}
\textbf{By the end of this chapter:} \emph{I can talk about a court, a lawyer, a sentence and prison, and say that someone was accused.}

Complete the last lesson of chapter 293: I can talk about a court, a lawyer, a sentence and prison, and say that someone was accused.
\end{chapteropening}

\section[\texorpdfstring{\ru{суд} (sud)}{суд (sud)}]{\ru{суд} (sud) — a court; a trial}

\label{lesson:RU-C293-sud}



\begin{quote}
Before the new one: say the Russian for questioning, then the Russian for a statement.
\end{quote}

\subsection*{You'll want to know: \ru{суд}}
\textbf{\ru{суд}} (\emph{sud}) — \textquotedblleft{}a court; a trial\textquotedblright{}.

\begin{quote}
\ru{Суд} \ru{будет} \ru{в} \ru{понедельник}.
\end{quote}

\begin{quote}
\emph{Sud búdet v ponedélnik.}
\end{quote}

\begin{quote}
The trial will be on Monday.
\end{quote}

\emph{\ru{В} \ru{суде}} — in court; \emph{\ru{подать} \ru{в} \ru{суд}} — to sue.

\begin{cousinweb}
\emph{\ru{суд}} is inherited through Old East Slavic \emph{sudŭ} from Proto-Slavic \emph{sǫdŭ}, which EtymDB builds from PIE \emph{som-} \textquotedblleft{}together\textquotedblright{} and \emph{d\textsuperscript{h}eh\textsubscript{1}-} \textquotedblleft{}to put\textquotedblright{}: a putting together.
\end{cousinweb}

\begin{grammarlens}[title={What you've built}]
One more word for telling, explaining and describing at B1.
\end{grammarlens}

\subsection*{Your turn}
\begin{itemize}
\raggedright
  \item \emph{Say it:} \emph{sud}
  \item \emph{Say it:} \emph{sud}, once more
  \item \emph{Say it:} a sentence of your own with \emph{sud}
  \item \emph{Recall:} say the Russian for questioning, then the Russian for a statement, then say \emph{sud} again
\end{itemize}

\begin{tcolorbox}[breakable,colback=teal!4,colframe=teal!35!black,title={Before you move on}]
What does \ru{суд} mean? (\textquotedblleft{}A court; a trial\textquotedblright{}.) Say it once more.
\end{tcolorbox}
\section[\texorpdfstring{\ru{адвокат} (advokát)}{адвокат (advokát)}]{\ru{адвокат} (advokát) — a lawyer (for the defence)}

\label{lesson:RU-C293-advokat}



\begin{quote}
Before the new one: say the Russian for a statement, then the Russian for a court.
\end{quote}

\subsection*{You'll want to know: \ru{адвокат}}
\textbf{\ru{адвокат}} (\emph{advokát}) — \textquotedblleft{}a lawyer (for the defence)\textquotedblright{}.

\begin{quote}
\ru{Адвокат} \ru{говорил} \ru{спокойно} \ru{и} \ru{уверенно}.
\end{quote}

\begin{quote}
\emph{Advokát govoríl spokóyno i uvérenno.}
\end{quote}

\begin{quote}
The lawyer spoke calmly and confidently.
\end{quote}

Masculine, also of a woman.

\begin{cousinweb}
\emph{\ru{адвокат}} goes back to Latin \emph{advocātus} (EtymDB), \textquotedblleft{}one called in to help\textquotedblright{}, the participle of \emph{advocāre} \textquotedblleft{}to call to\textquotedblright{} (Lewis \& Short).
\end{cousinweb}

\begin{grammarlens}[title={What you've built}]
One more word for telling, explaining and describing at B1.
\end{grammarlens}

\subsection*{Your turn}
\begin{itemize}
\raggedright
  \item \emph{Say it:} \emph{advokát}
  \item \emph{Say it:} \emph{advokát}, once more
  \item \emph{Say it:} a sentence of your own with \emph{advokát}
  \item \emph{Recall:} say the Russian for a statement, then the Russian for a court, then say \emph{advokát} again
\end{itemize}

\begin{tcolorbox}[breakable,colback=teal!4,colframe=teal!35!black,title={Before you move on}]
What does \ru{адвокат} mean? (\textquotedblleft{}A lawyer (for the defence)\textquotedblright{}.) Say it once more.
\end{tcolorbox}
\section[\texorpdfstring{\ru{приговор} (prigovór)}{приговор (prigovór)}]{\ru{приговор} (prigovór) — a sentence, a verdict}

\label{lesson:RU-C293-prigovor}



\begin{quote}
Before the new one: say the Russian for a court, then the Russian for a lawyer.
\end{quote}

\subsection*{You'll want to know: \ru{приговор}}
\textbf{\ru{приговор}} (\emph{prigovór}) — \textquotedblleft{}a sentence, a verdict\textquotedblright{}.

\begin{quote}
\ru{Приговор} \ru{был} \ru{очень} \ru{строгим}.
\end{quote}

\begin{quote}
\emph{Prigovór byl óchen strógim.}
\end{quote}

\begin{quote}
The sentence was very harsh.
\end{quote}

\emph{\ru{Вынести} \ru{приговор}} — to pass sentence.

\begin{cousinweb}
\emph{\ru{приговор}} is built on \emph{\ru{при}-} and the root of \emph{\ru{говорить}} \textquotedblleft{}to speak\textquotedblright{}, inherited from Proto-Slavic \emph{govoriti} (EtymDB): what is pronounced.
\end{cousinweb}

\begin{grammarlens}[title={What you've built}]
One more word for telling, explaining and describing at B1.
\end{grammarlens}

\subsection*{Your turn}
\begin{itemize}
\raggedright
  \item \emph{Say it:} \emph{prigovór}
  \item \emph{Say it:} \emph{prigovór}, once more
  \item \emph{Say it:} a sentence of your own with \emph{prigovór}
  \item \emph{Recall:} say the Russian for a court, then the Russian for a lawyer, then say \emph{prigovór} again
\end{itemize}

\begin{tcolorbox}[breakable,colback=teal!4,colframe=teal!35!black,title={Before you move on}]
What does \ru{приговор} mean? (\textquotedblleft{}A sentence, a verdict\textquotedblright{}.) Say it once more.
\end{tcolorbox}
\section[\texorpdfstring{\ru{тюрьма} (tyurmá)}{тюрьма (tyurmá)}]{\ru{тюрьма} (tyurmá) — a prison}

\label{lesson:RU-C293-tyurma}



\begin{quote}
Before the new one: say the Russian for a lawyer, then the Russian for a sentence.
\end{quote}

\subsection*{You'll want to know: \ru{тюрьма}}
\textbf{\ru{тюрьма}} (\emph{tyurmá}) — \textquotedblleft{}a prison\textquotedblright{}.

\begin{quote}
\ru{Он} \ru{провёл} \ru{в} \ru{тюрьме} \ru{пять} \ru{лет}.
\end{quote}

\begin{quote}
\emph{On provyól v tyurmé pyat let.}
\end{quote}

\begin{quote}
He spent five years in prison.
\end{quote}

Genitive plural \emph{\ru{тюрем}}.

\begin{cousinweb}
\emph{\ru{тюрьма}} is traced by EtymDB to Proto-Turkic \emph{türmä} \textquotedblleft{}prison\textquotedblright{}; others derive it through Polish from German \emph{Turm} \textquotedblleft{}a tower\textquotedblright{}, so the origin is not fully settled.
\end{cousinweb}

\begin{grammarlens}[title={What you've built}]
One more word for telling, explaining and describing at B1.
\end{grammarlens}

\subsection*{Your turn}
\begin{itemize}
\raggedright
  \item \emph{Say it:} \emph{tyurmá}
  \item \emph{Say it:} \emph{tyurmá}, once more
  \item \emph{Say it:} a sentence of your own with \emph{tyurmá}
  \item \emph{Recall:} say the Russian for a lawyer, then the Russian for a sentence, then say \emph{tyurmá} again
\end{itemize}

\begin{tcolorbox}[breakable,colback=teal!4,colframe=teal!35!black,title={Before you move on}]
What does \ru{тюрьма} mean? (\textquotedblleft{}A prison\textquotedblright{}.) Say it once more.
\end{tcolorbox}
\section[\texorpdfstring{\ru{обвинять} (obvinyát)}{обвинять (obvinyát)}]{\ru{обвинять} (obvinyát) — to accuse}

\label{lesson:RU-C293-obvinyat}



\begin{quote}
Before the new one: say the Russian for a sentence, then the Russian for a prison.
\end{quote}

\subsection*{You'll want to know: \ru{обвинять}}
\textbf{\ru{обвинять}} (\emph{obvinyát}) — \textquotedblleft{}to accuse\textquotedblright{}.

\begin{quote}
\ru{Его} \ru{обвиняют} \ru{в} \ru{краже}.
\end{quote}

\begin{quote}
\emph{Yevó obvinyáyut v krázhe.}
\end{quote}

\begin{quote}
He is accused of theft.
\end{quote}

Imperfective; perfective \emph{\ru{обвинить}} (\emph{obvinít}). The person is in the accusative, the charge with \emph{\ru{в}} + prepositional.

\begin{cousinweb}
\emph{\ru{обвинять}} is built on \emph{\ru{об}-} and \emph{\ru{вина}} \textquotedblleft{}guilt\textquotedblright{}, an inherited Slavic word: to lay guilt on someone. EtymDB lists no parent for the verb.
\end{cousinweb}

\begin{grammarlens}[title={What you've built}]
One more word for telling, explaining and describing at B1.
\end{grammarlens}

\subsection*{Your turn}
\begin{itemize}
\raggedright
  \item \emph{Say it:} \emph{obvinyát}
  \item \emph{Say it:} \emph{obvinyát}, once more
  \item \emph{Say it:} a sentence of your own with \emph{obvinyát}
  \item \emph{Recall:} say the Russian for a sentence, then the Russian for a prison, then say \emph{obvinyát} again
\end{itemize}

\begin{tcolorbox}[breakable,colback=teal!4,colframe=teal!35!black,title={Before you move on}]
What does \ru{обвинять} mean? (\textquotedblleft{}To accuse\textquotedblright{}.) Say it once more.
\end{tcolorbox}
